\chapter{总结与展望}\label{chap:conclusions}
\section{工作总结}
本文围绕一套合成无人机时序数据组织了完整的排版案例。
在数据层面，按照架次划分训练、校准和测试，避免跨边界滑窗，
并将训练统计量用于全部子集的标准化。
在方法层面，建立了岭回归预测、残差聚合、因果平滑和分位数阈值的透明流程。
在结果层面，所有表格和曲线由同一份脚本生成，关键数值通过宏进入摘要和正文。

从模板角度看，案例覆盖了多级标题、双语图表题注、子图引用、三线表、
公式与矩阵、定理证明、算法、代码、跨页长表和两套文献后端。
这些接口的共同作用是使内容调整与版式控制相互分离。
更换实验数据后，用户需要更新数据生成或读取逻辑，而不应重新手工编号图表。

\section{尚未解决的问题}
合成数据实验只验证程序链路是否一致，并不验证真实飞行故障检测的可靠性。
在开展实际研究时，需要重新定义数据来源、标注协议、平台差异和异常类别，
并记录可复现的数据划分与预处理步骤。
评价还应覆盖与应用有关的报警延迟、误报负担及不同架次上的稳定性。

模板本身仍需要对照最终提交要求检查所有页面。
本示例未提供原创性声明、授权声明或答辩委员会的正式文本。
中文与英文摘要也由示例辅助命令实现，后续可迁移为正式环境。
字体、页边距、题注语言和编号配置是否满足要求，应由具体版本的学校规范确认。

\section{后续研究与工程方向}
研究层面可以引入真实飞行记录、分析正常机动与异常的区分，
并在一致协议下比较不同预测和重构模型。
工程层面应建立持续的编译回归测试，检查四种学位组合、两套文献后端、
中文与双语题注以及附录编号是否保持一致。
对于文件组织，宜将数据、配置、算法实现和排版源码分别维护，
再通过构建脚本连接，减少难以追踪的手工修改。
